For SQ2, the specified feedback rule changes exposure without producing an advantage in exposure share relative to production in the one-item-feed grid. The clearest high-activity example shows why the attempt measurements matter: posting acceptance and receipts per attempt increase, while receipts per accepted post decrease. The extra accepted output can raise the subgroup's receipt share even when \(D\) falls. We therefore cannot treat a negative \(D\) as evidence that the producer receives less attention in every sense, or treat higher receipt share as evidence of disproportionate amplification.

The adaptive rule is operationally different from high activity alone. It uses past receipt outcomes, and the implementation tests show that changing those outcomes changes future bin choices. The observed concentration of individual choices is consistent with strong exploitation of earlier successes. However, responsiveness to feedback does not guarantee improved performance. Delayed cascades, sparse exploration, changing network conditions and age differences among posts can make a cumulative per-post score favour choices that do not improve overall exposure efficiency.

For SQ3, we do not resolve an interaction between policy and activity over the tested \(2\times\)--\(10\times\) contrast. We also do not resolve the adaptive effect at capacities 3 and 5. These uncertain estimates limit transfer of the one-item-feed comparisons; they do not prove that capacity has removed an effect or that the policy is equivalent to its controls. The separate activity-only capacity experiment further shows that a simple limited-attention explanation is insufficient to identify the underlying mechanism.

